problem:  
Let \((X,d,\mu)\) be a complete, proper, **Ahlfors \(N\)-regular** geodesic metric measure space satisfying the **Bishop–Gromov volume comparison inequality** with parameters \((K,N)\):
\[
r \mapsto \frac{\mu(B(x,r))}{v_{K,N}(r)} \quad \text{is nonincreasing for all } x\in X.
\]
Assume further that tangent measures \(\mathsf{Tan}(\mu,x)\) exist at \(\mu\)-almost every \(x\in X\).  
Define the **regular set**
\[
\mathcal{R} := \Bigl\{ x\in X : \text{there exists a tangent measure of the form } c_x\,\mathcal{H}^N\!\lfloor_{\mathbb{R}^N}\Bigr\},
\]
and singular set \(\mathcal{S} := X\setminus \mathcal{R}\).

**Open problem:**  
Does the Bishop–Gromov volume comparison property alone imply that the singular set \(\mathcal{S}\) has **Minkowski codimension at least 1**, i.e. there exists \(\eta>0\) such that
\[
\limsup_{r\to 0} \frac{\mu(B(\mathcal{S},r))}{r^{N-1-\eta}} = 0 ?
\]
Equivalently: can one construct an Ahlfors \(N\)-regular metric measure space obeying Bishop–Gromov comparison but whose singular set has **positive \(N\)-dimensional density** or Minkowski dimension \(N\)?  

Why is it a "good" problem:  
This question isolates a sharp, conceptual boundary of what *pure volume comparison* can enforce about the singular geometry of non‑smooth spaces.

- **Profound insight:** In Ricci limit and \(RCD(K,N)\) spaces, analytic curvature conditions imply strong stratification (codimension ≥ 2) and rectifiability of regular parts. The problem asks whether removing all analytic assumptions—keeping only the Bishop–Gromov monotonicity—still forces any dimensional gap at all. That is, does the coarse global constraint on volume growth already prevent the presence of “thick” singular sets? Either answer would reinterpret the Bishop–Gromov principle as a genuinely geometric regulator or expose its precise limitations.  

- **Bridge between fields:** Addressing the problem links **comparison geometry**, **Preiss’s tangent‑measure theory**, and **geometric measure theory** on metric spaces. It draws techniques from analytic curvature bounds, quantitative stratification, and fractal geometry, creating a natural meeting point between these areas.  

- **Extreme simplicity:** The statement refers only to balls, measures, tangent scaling, and limiting dimension—concepts accessible at a basic geometric level—yet its resolution would demand new structural theorems tying synthetic curvature comparison to measure‑theoretic dimensional decay. In this sense, it captures the ideal of “narrow entrance, profound depth.”